\documentclass[11pt]{amsart}
\usepackage{amsmath,amssymb,amsthm}
\usepackage[margin=1.05in]{geometry}
\raggedbottom

\numberwithin{equation}{section}
\newtheorem{theorem}{Theorem}[section]
\newtheorem{proposition}[theorem]{Proposition}
\newtheorem{lemma}[theorem]{Lemma}
\newtheorem{corollary}[theorem]{Corollary}
\theoremstyle{definition}
\newtheorem{definition}[theorem]{Definition}
\theoremstyle{remark}
\newtheorem{remark}[theorem]{Remark}
\newtheorem{example}[theorem]{Example}

\newcommand{\Z}{\mathbb{Z}}
\newcommand{\Q}{\mathbb{Q}}
\newcommand{\R}{\mathbb{R}}
\newcommand{\N}{\mathbb{N}}
\newcommand{\C}{\mathbb{C}}
\newcommand{\Cp}{\mathbb{C}_p}
\newcommand{\cA}{\mathcal{A}}
\newcommand{\cO}{\mathcal{O}}
\newcommand{\pp}{\mathfrak{p}}
\newcommand{\qq}{\mathfrak{q}}
\newcommand{\aaa}{\mathfrak{a}}
\newcommand{\Ns}{\mathrm{N}}
\newcommand{\Cl}{\operatorname{Cl}}
\newcommand{\Gal}{\operatorname{Gal}}
\newcommand{\Frob}{\operatorname{Frob}}
\newcommand{\Hom}{\operatorname{Hom}}
\newcommand{\rank}{\operatorname{rank}}
\newcommand{\chr}{\operatorname{char}}
\newcommand{\id}{\mathrm{id}}

\PassOptionsToPackage{hyphens}{url}
\usepackage[hidelinks,bookmarks=false]{hyperref}

\title[No $p$-adic KMS states]{Iwasawa theory and localized Bost--Connes systems:\\ there are no $p$-adic KMS states,\\ and where Leopoldt actually enters}
\author{Ruqing Chen}
\address{GUT Geoservice Inc., Montreal, Canada}
\email{ruqing@hotmail.com}
\thanks{Sources, code and data for the entire series: \url{https://github.com/Ruqing1963/localized-bost-connes}; this paper is in \texttt{papers/XXII-iwasawa/}.}
\date{}

\begin{document}

\begin{abstract}
We examine the proposal to build a $p$-adic thermodynamics on the localized Bost--Connes
systems and to read the Iwasawa invariants off it. Four obstructions close it, of which the
first is fundamental and the others are independent of it.

\emph{There are no $p$-adic KMS states.} The KMS condition is a statement about a
\emph{state} --- a positive normalized functional --- on a $C^*$-algebra, together with the
analytic continuation of $\sigma_t$ to a complex strip. Positivity requires an ordered field
and an involution compatible with it; $\Cp$ is not ordered and carries no $C^*$-theory. There
is $p$-adic functional analysis, but no positive cone, hence no states, hence no simplex, and
hence nothing on which a KMS condition could be imposed.

\emph{The proposed dynamics is undefined on the proposed set.} The map $x\mapsto\langle x\rangle$
projects $\Z_p^\times$ onto $1+p\Z_p$ and requires $x$ to be a $p$-adic unit. For
$\pp\mid p$ one has $\Ns\pp=p^{f}$ with $v_p(\Ns\pp)=f>0$, so $\langle\Ns\pp\rangle^z$ does
not exist. The localization $S=\{\pp\mid p\}$ is exactly the set at which the $p$-adic
dynamics cannot be defined --- which is why $p$-adic $L$-functions remove the Euler factors
at $p$.

\emph{There is no $p$-adic abscissa of convergence.} For $\qq\nmid p$ the number
$\langle\Ns\qq\rangle^{-z}$ is a $1$-unit, so its absolute value is $1$ for every $z$;
a $p$-adic series converges iff its terms tend to $0$, and terms of absolute value $1$ never
do. Every $p$-adic Dirichlet series over an infinite set of primes diverges, at every weight.
The convergence--divergence dichotomy that drives the whole archimedean theory therefore has
no $p$-adic counterpart. This is not a defect of our set-up: it is why $p$-adic
$L$-functions are constructed from measures and Iwasawa power series rather than from
Dirichlet series.

\emph{The proposed obstruction cannot generate the characteristic ideal.} The expression
\[
\Xi_p(z,\chi)=\sum_{\pp\mid p}\bigl(1-\chi(\Frob_\pp)\langle\Ns\pp\rangle^z\bigr)
\]
is a finite sum of $g\le[K:\Q]$ terms, so its $\lambda$-invariant is bounded in terms of $g$
alone; the Iwasawa
$\lambda$-invariant is unbounded. A bounded-degree element generates no family of ideals of
unbounded degree.

What survives is smaller and precise. The $p$-local partition function
$\prod_{\pp\mid p}(1-\Ns\pp^{-\beta})^{-1}$ is, up to inversion, exactly the Euler factor that
$L_p(s,\chi)$ removes in its interpolation formula: a genuine dictionary entry, but between
the partition function and the interpolation factor, not the characteristic ideal. And
Leopoldt's conjecture does enter this programme --- not through a $p$-adic state space but
through Proposition 2.2 and Remark 2.3 of Paper II of the series, where the isotropy of the
groupoid at the boundary strata
requires a Leopoldt-type condition for general $K$ and is unconditional when $\cO_K^\times$
is finite.
\end{abstract}

\maketitle

\section{There are no $p$-adic KMS states}

Papers II, III, IV, VIII, X and XXI of the series, referred to by numeral, are unpublished
companions; every fact used from them is restated or verified in place. \cite{Main} and
Paper I \cite{SeriesI} are cited by DOI.

\begin{proposition}[Positivity is unavailable]\label{prop:nostates}
Let $A$ be a Banach algebra over $\Cp$ with an involution. There is no notion of state on
$A$: a state is a linear functional $\varphi$ with $\varphi(x^*x)\ge0$ and $\varphi(1)=1$, and
the inequality requires an order on the scalar field. The field $\Cp$ is not orderable: it is
algebraically closed, so $-1=i^2$ is a square; and already $\Q_p$ is non-orderable, $-1$
being a sum of at most four squares (the level of $\Q_p$ is $1$, $2$ or $4$ according to
$p\equiv1\ (4)$, $p\equiv3\ (4)$, $p=2$). Consequently there is no positive cone, no state
space, no Choquet simplex, and no
Gelfand--Naimark--Segal construction.
\end{proposition}

\begin{corollary}[The KMS condition has no $p$-adic form]\label{cor:noKMS}
The $\mathrm{KMS}_\beta$ condition presupposes (i) a state $\varphi$ satisfying
$\varphi(xy)=\varphi(y\,\sigma_{i\beta}(x))$ and (ii) the extension of $t\mapsto\sigma_t$ to
the strip $0\le\Im t\le\beta$ in $\C$. Neither is
available over $\Cp$: (i) by Proposition~\ref{prop:nostates}, and (ii) because $\Cp$ has no
distinguished strip, the ``imaginary direction'' being an artefact of $\C=\R\oplus i\R$ with
no $p$-adic analogue. There is therefore no object called a $p$-adic KMS state, and no
$p$-adic KMS simplex whose dimension could be computed.
\end{corollary}

\begin{remark}[What can be transported]\label{rem:transport}
The Gibbs functional $\varphi_\beta(e_\aaa)=\Ns\aaa^{-\beta}$ does have a $p$-adic avatar,
but as a \emph{measure} rather than a state: the assignment
$\aaa\mapsto\langle\Ns\aaa\rangle^{-z}$ extends to a $\Cp$-valued distribution on a compact
space, and it is exactly such distributions --- the Mazur measure and its relatives --- that
define $p$-adic $L$-functions. Measures do not form a simplex without positivity, and no
equilibrium condition selects them. The correct statement of the analogy is therefore
``Gibbs state $\leftrightarrow$ $p$-adic measure'', with the equilibrium characterization lost
in transit.
\end{remark}

\section{The dynamics is undefined on $S=\{\pp\mid p\}$}

\begin{proposition}[$\langle\Ns\pp\rangle$ does not exist]\label{prop:undefined}
The projection $\langle\,\cdot\,\rangle:\Z_p^\times\to1+p\Z_p$, $x\mapsto x/\omega(x)$, is
defined only on $p$-adic units. For $\pp\mid p$ with residue degree $f$ one has
$\Ns\pp=p^{f}$ and $v_p(\Ns\pp)=f>0$, so $\Ns\pp\notin\Z_p^\times$ and
$\langle\Ns\pp\rangle^z$ is undefined. Hence the proposed $p$-adic dynamics
$\sigma^{(p)}_z(\mu_\aaa)=\langle\Ns\aaa\rangle^z\mu_\aaa$ does not exist on
$J_p=\bigoplus_{\pp\mid p}\N\pp$.
\end{proposition}

\begin{remark}[The localization is at the wrong set]\label{rem:wrongset}
This is not a normalization issue to be repaired by a choice of branch. The set
$\{\pp\mid p\}$ is precisely the set of places at which $p$-adic interpolation breaks down,
which is why $L_p(s,\chi)$ is defined with the Euler factors at $p$ \emph{removed}
\cite{Washington}:
\[
L_p(1-n,\chi\omega^{-n})=L(1-n,\chi\omega^{-n})\prod_{\pp\mid p}\bigl(1-\chi\omega^{-n}(\pp)\Ns\pp^{\,n-1}\bigr).
\]
A $p$-adic dynamics must be built on the primes $\qq\nmid p$, where $\Ns\qq\in\Z_p^\times$.
That, however, runs into \S\ref{sec:noabscissa}.
\end{remark}

\section{There is no $p$-adic abscissa of convergence}\label{sec:noabscissa}

\begin{theorem}[$p$-adic Dirichlet series never converge]\label{thm:noabscissa}
Let $T$ be an infinite set of primes with $\qq\nmid p$ for all $\qq\in T$, and let
$z$ range over the domain on which $u^z=\exp_p(z\log_pu)$ is defined for $1$-units $u$ ---
all of $\Z_p$, and the disc $|z|_p<p^{(p-2)/(p-1)}$ of $\Cp$. Then
$\langle\Ns\qq\rangle\in1+p\Z_p$, so
\[
\bigl|\langle\Ns\qq\rangle^{-z}\bigr|_p=1\qquad\text{for every }\qq\in T\text{ and every such }z,
\]
and since a series in a complete non-archimedean field converges if and only if its terms
tend to $0$, the series $\sum_{\qq\in T}\langle\Ns\qq\rangle^{-z}$ diverges for every $z$.
The same holds for any weighted series
$\sum_{\qq\in T}w_\qq\,\langle\Ns\qq\rangle^{-z}$ with $|w_\qq|_p=1$ --- for instance
$w_\qq=1-\chi(\Frob_\qq)$ for a character $\chi$ with $\chi(\Frob_\qq)\ne1$ of order prime
to $p$.
\end{theorem}

\begin{corollary}[No $p$-adic phase transition]\label{cor:nophase}
There is no $p$-adic critical exponent, no convergence--divergence dichotomy, and hence no
$p$-adic Kakutani criterion. The mechanism that produces every phase transition in
\cite{Main}, \cite{SeriesI} and Paper IV of the series --- the competition between the decay of
$\Ns\pp^{-\beta}$ and the abundance of primes --- is archimedean, and it is absent
$p$-adically because all the terms have absolute value $1$.
\end{corollary}

\begin{remark}[Why $p$-adic $L$-functions look the way they do]\label{rem:whymeasures}
Corollary~\ref{cor:nophase} is not a peculiarity of Bost--Connes systems. It is the reason
$p$-adic $L$-functions are never defined by Dirichlet series: Kubota--Leopoldt, Iwasawa
\cite{Iwasawa} and
Mazur all construct them from measures on $\Z_p^\times$, or equivalently from power series in
$\Lambda=\Z_p[[T]]$, precisely because the naive series has no region of convergence. Any
$p$-adic thermodynamics would have to be built on measures, and would then lack the
variational characterization that makes the archimedean theory informative.
\end{remark}

\section{The obstruction cannot generate the characteristic ideal}

\begin{proposition}[The element cannot generate the characteristic ideal]\label{prop:bounded}
Suppose, counterfactually to Proposition~\ref{prop:undefined}, that
$\Xi_p(z,\chi)=\sum_{i=1}^g(1-\chi_i\,u_i^z)$ were defined, with $u_i$ $1$-units,
$\chi_i=\chi(\Frob_{\pp_i})$, and $g=\#\{\pp\mid p\}\le[K:\Q]$; under $u_i^z=(1+T)^{c_i}$ it
becomes an element of $\Lambda\otimes\Z_p[\chi]$. It cannot generate the characteristic
ideal $\chr_\Lambda(X)=(p^{\mu}P(T))$ of the main conjecture \cite{MazurWiles}, for two
independent reasons.
\begin{enumerate}
\item Already for $K=\Q$, $g=1$: if $\chi(p)\ne1$ has order prime to $p$ then
$|1-\chi(p)(1+T)^{c}|_p=|1-\chi(p)|_p=1$ on the open unit disc, so the element is a
\emph{unit} of $\Lambda\otimes\Z_p[\chi]$ and generates the unit ideal --- while
$\chr_\Lambda(X_\chi)$ is a proper ideal whenever $\lambda_\chi+\mu_\chi>0$, which occurs.
\item In general the element depends only on the $g$ pairs
$(\chi(\Frob_{\pp_i}),\langle\Ns\pp_i\rangle)$ --- finitely many $p$-adic parameters
attached to the decomposition of $p$ in $K$ --- whereas the characteristic ideal encodes the
Iwasawa module of the full tower, which is not a function of the decomposition data at $p$:
fields with identical splitting behaviour at $p$ have different $\lambda$-invariants.
\end{enumerate}
We note that the tempting quantitative version --- ``a sum of $g$ binomials has
$\lambda$-invariant bounded in terms of $g$'' --- is itself false: already
$(1+T)^{c'}-(1+T)^{c}=(1+T)^{c}\bigl((1+T)^{c'-c}-1\bigr)$ has
$\lambda=p^{v_p(c'-c)}$, unbounded; the failure of the proposal is the wrong
\emph{dependence}, not a degree count.
\end{proposition}

\begin{remark}[And $\mu$]\label{rem:mu}
The proposed formula $\mu_p(K)=v_p(\mathrm{Res}_{z=0}\zeta_{p,\mathrm{KMS}}(z))$ has no
content, there being no $\zeta_{p,\mathrm{KMS}}$ by Corollary~\ref{cor:noKMS}. We note also
that for abelian $K/\Q$ the Ferrero--Washington theorem \cite{FW} gives $\mu_p(K)=0$, so any
correct
formula must return $0$ in that case; a residue at a pole of a nonexistent function does not.
The genuine residue statement of $p$-adic $L$-theory is Colmez's formula for the residue of
the $p$-adic zeta function of $K$ at $s=1$ \cite{Colmez}, and what appears there is the
$p$-adic regulator $R_p(K)$ --- the Leopoldt quantity of \S\ref{sec:true} --- not $\mu$.
\end{remark}

\section{What is true}\label{sec:true}

\begin{proposition}[The partition function is the interpolation factor]\label{prop:euler}
For $S=\{\pp\mid p\}$ the partition function of the localized system is the finite product
\[
\zeta_{K,p}(\beta)=\prod_{\pp\mid p}\bigl(1-\Ns\pp^{-\beta}\bigr)^{-1},
\]
whose inverse, twisted by $\chi$, is exactly the Euler factor removed in the interpolation
formula for $L_p(s,\chi)$ recalled in Remark~\ref{rem:wrongset}. This is a precise
dictionary entry: \emph{$p$-local Bost--Connes partition function $\leftrightarrow$ $p$-adic
interpolation factor.} It is not a dictionary entry for the characteristic ideal.
\end{proposition}

\begin{remark}[Finite $S$ has no thermodynamics]\label{rem:finiteS}
Consistently with this, $S=\{\pp\mid p\}$ is finite, so $\beta_c=0$ and by
the finite-generation exclusion of Paper X (its Remark 1.5) --- a finite generating set gives a finite obstruction sum ---
$\Xi_\beta=\widehat{G_p}$ for every $\beta>0$ and the chain is constant. The archimedean
$p$-local system has no phase transition, which is what motivated looking for a $p$-adic one;
Corollary~\ref{cor:nophase} says the search cannot succeed in that direction.
\end{remark}

\begin{proposition}[Where Leopoldt genuinely enters]\label{prop:leopoldt}
Let $M_\infty$ be the maximal abelian pro-$p$ extension of $K$ unramified outside $p$. Then
\[
\rank_{\Z_p}\Gal(M_\infty/K)=r_2+1+\delta_p(K),
\]
with $\delta_p(K)\ge0$ the Leopoldt defect, and Leopoldt's conjecture is $\delta_p(K)=0$,
equivalently $R_p(K)\ne0$. This quantity does control the localized Bost--Connes system, but
in the archimedean theory: by Proposition 2.2 and Remark 2.3 of Paper II the isotropy of $\mathcal G_S$ off
the boundary strata is computed from the closure $\overline{\cO^\times_{K,+}}$ inside
$U_S$, and the computation requires a Leopoldt-type condition for general $K$, being
unconditional when $\cO_K^\times$ is finite, that is for $K=\Q$ and imaginary quadratic $K$.
\end{proposition}

\begin{remark}\label{rem:leopoldtplace}
So the intuition behind the proposal is sound --- Leopoldt does belong to this story --- but
its place is the isotropy computation of Paper II, not the dimension of a $p$-adic
state space. That the topological principality of $\mathcal G_S$ was established there
\emph{independently} of Leopoldt (its Theorem 2.5), the strata being nowhere dense regardless, is what keeps the
main classification unconditional.
\end{remark}

\section{Assessment}

\[
\begin{array}{lll}
\text{Proposal} & \text{Status} & \text{Reference}\\\hline
p\text{-adic KMS states} & \textbf{do not exist: no positivity} & \text{Prop.~1.1, Cor.~1.2}\\
\sigma^{(p)}_z\ \text{on}\ \{\pp\mid p\} & \textbf{undefined: }\Ns\pp\notin\Z_p^\times & \text{Prop.~2.1}\\
p\text{-adic Dirichlet series} & \textbf{never converge} & \text{Thm.~3.1}\\
p\text{-adic }\beta_c,\ \text{phase transition} & \textbf{do not exist} & \text{Cor.~3.2}\\
(\Xi_p)=\chr_\Lambda(X) & \textbf{false: wrong dependence} & \text{Prop.~4.1}\\
\mu\ \text{from a residue} & \text{no content} & \text{Rem.~4.2}\\
\lambda\ \text{as a count of transitions} & \text{no content} & \text{Cor.~3.2}\\
\rank\Gal(M_\infty/K)=r_2+1+\delta_p & \text{true} & \text{Prop.~5.3}\\[-1pt]
\text{Leopoldt controls the system} & \textbf{true, but via Paper II} & \text{Prop.~5.3}\\
\zeta_{K,p}\leftrightarrow\text{interpolation factor} & \textbf{true} & \text{Prop.~5.1}
\end{array}
\]

The proposal fails at a level below its own statements. Everything in the archimedean theory
--- KMS states, simplices, phase transitions, obstruction groups --- rests on positivity and
on the convergence--divergence dichotomy of real Dirichlet series, and neither has a $p$-adic
form. The absence is not a gap in the technology but a fact about $\Cp$: it is not ordered,
and its units all have absolute value $1$.

That said, the underlying instinct identifies two real connections, and both are worth
recording in the form they actually take. The $p$-local partition function is the
interpolation factor of $L_p$, which is a statement about Euler factors and is exact.
And Leopoldt's conjecture does govern part of this programme, namely the isotropy computation
of Paper II (its Proposition 2.2 and Remark 2.3), which is where the units enter --- consistently with
Paper VIII, whose inertia model, built on the meridian direction matched to the global
units, was shown to carry no thermodynamics (its Theorem 6.1 and the moral of its
Remark 6.3), and with Paper X (its Remark 1.5), where the finiteness of the unit group
decided which semigroups are admissible. The units recur; the $p$-adic states do not exist.

Two questions we would consider worth asking instead. Is there a $p$-adic analogue of the
\emph{measure class} invariant of Corollary 2.4 of Paper III --- since that, unlike a state,
is not a positivity notion, and since the whole corpus has found the arithmetic to live
there? And second, the archimedean $p$-local system is trivial by
Remark~\ref{rem:finiteS}; is there an intermediate localization, $S$ infinite but with the
$p$-adic weights attached as a coefficient system rather than as the dynamics, in which both
structures coexist? That is the closing suggestion of Paper XXI (its \S5), itself pointing
at \S6 of Paper VI, and for the same
reason: the arithmetic input belongs in the coefficients, not in the flow.

\begin{thebibliography}{9}
\bibitem{Colmez} P. Colmez, \emph{R\'esidu en $s=1$ des fonctions z\^eta $p$-adiques}, Invent. Math. 91 (1988), 371--389.
\bibitem{FW} B. Ferrero, L. Washington, \emph{The Iwasawa invariant $\mu_p$ vanishes for abelian number fields}, Ann. of Math. (2) 109 (1979), 377--395.
\bibitem{Iwasawa} K. Iwasawa, \emph{On $\Z_\ell$-extensions of algebraic number fields}, Ann. of Math. (2) 98 (1973), 246--326.
\bibitem{MazurWiles} B. Mazur, A. Wiles, \emph{Class fields of abelian extensions of $\Q$}, Invent. Math. 76 (1984), 179--330.
\bibitem{Washington} L. C. Washington, \emph{Introduction to Cyclotomic Fields}, 2nd ed., Graduate Texts in Math. 83, Springer, 1997.
\bibitem{Main} R. Chen, \emph{KMS state uniqueness and phase transitions in localized Bost--Connes systems}, preprint, 2026, \newline\url{https://doi.org/10.5281/zenodo.22152101}.
\bibitem{SeriesI} R. Chen, \emph{Localized Bost--Connes systems I: factor type and the spectrum of transitions}, preprint, 2026, DOI \href{https://doi.org/10.5281/zenodo.22160827}{10.5281/zenodo.22160827}.
\end{thebibliography}

\end{document}
